\chapter*{Kết luận}
\addcontentsline{toc}{chapter}{Kết luận}

\section*{1. Kết quả đạt được}

Đồ án \textbf{\projecttitle} đã giúp nhóm vận dụng quy trình công nghệ phần mềm vào một bài toán quản lý thực tế. Nhóm đã thực hiện từ khảo sát yêu cầu, chốt phạm vi, thiết kế hệ thống, thiết kế dữ liệu, thiết kế xử lý, thiết kế giao diện đến hiện thực, kiểm thử và triển khai local.

Kết quả kỹ thuật đạt được gồm:

\begin{itemize}
    \item Hệ thống web nội bộ hỗ trợ 8 vai trò người dùng: ADMIN, RECEPTIONIST, DOCTOR, CASHIER, MANAGER, NURSE, LAB\_TECHNICIAN, PHARMACIST.
    \item Backend NestJS được tổ chức theo modular monolith với 21 feature folders, 20 controllers có API route và 21 service files.
    \item Database PostgreSQL được định nghĩa bằng Prisma với 37 models, 12 enums và 3 migrations.
    \item API có 92 endpoints, gồm 41 endpoints Phase 1 và 51 endpoints Phase 2, với prefix \code{/api/v1} và Swagger UI tại \code{/api/docs}.
    \item Frontend React có 36 named routes và 33 page files, tổ chức theo feature, có ProtectedRoute và RequireRole.
    \item Có 3 e2e test files có ý nghĩa cho các luồng lõi Phase 1; theo log nhóm cung cấp, test run tự động đạt 220/220 pass và cần được lưu bằng screenshot/log khi nộp.
\end{itemize}

\section*{2. Đánh giá so với mục tiêu ban đầu}

So với mục tiêu ban đầu, hệ thống đã đáp ứng được các nghiệp vụ lõi của phòng mạch tư nhân như đăng nhập, phân quyền, quản lý bệnh nhân, tiếp nhận lượt khám, lập phiếu khám, kê đơn, hóa đơn, thanh toán và báo cáo. Ngoài ra, Phase 2 mở rộng thêm các nghiệp vụ nâng cao như lịch hẹn, hàng đợi, sinh hiệu, dịch vụ, xét nghiệm, kho thuốc, cấp phát thuốc, tổ chức phòng khám và audit log. Xét theo phạm vi yêu cầu, báo cáo đã truy vết được 30 yêu cầu chức năng, 8 vai trò người dùng, 15 business rules trọng tâm và 92 endpoints đã hiện thực.

Việc chia hệ thống thành hai phase giúp nhóm kiểm soát phạm vi tốt hơn: Phase 1 đóng vai trò nền tảng nghiệp vụ lõi, Phase 2 mở rộng dựa trên nền tảng đã có. Cách tổ chức này phù hợp với quy trình phát triển phần mềm có kiểm soát.

\section*{3. Hạn chế}

Hệ thống vẫn còn một số hạn chế cần trình bày trung thực:

\begin{itemize}
    \item Khách hàng trong phạm vi đồ án là khách hàng giả định, không có phỏng vấn thực tế.
    \item Kiểm thử tự động chủ yếu mới bao phủ Phase 1; Phase 2 cần manual test evidence hoặc automated tests trong tương lai.
    \item Chưa có unit test đầy đủ cho service layer.
    \item Chưa có frontend automated tests.
    \item Chưa có Docker, docker-compose, CI/CD pipeline hoặc production deployment.
    \item Bằng chứng đóng góp nhóm được trình bày qua cả Git history, biên bản phân công, checklist kiểm thử, screenshot demo và tài liệu; không chỉ dựa vào commit count vì nhóm có thực hiện pair programming và các hoạt động ngoài commit code.
    \item Business rule hoàn tất phiếu khám hiện chưa chặn trường hợp còn pending service/lab orders; đây là limitation cần cải thiện nếu hệ thống vận hành thực tế.
\end{itemize}

\section*{4. Hướng phát triển}

Trong tương lai, nhóm có thể tiếp tục phát triển hệ thống theo các hướng:

\begin{itemize}
    \item Cổng bệnh nhân để bệnh nhân đặt lịch, xem lịch sử và nhận thông báo.
    \item Nhắc lịch qua SMS/email.
    \item Tích hợp bảo hiểm y tế hoặc bảo hiểm tư nhân.
    \item Multi-branch/multi-tenant cho nhiều chi nhánh phòng khám.
    \item Dashboard phân tích nâng cao và báo cáo BI.
    \item Bổ sung Docker, CI/CD pipeline, monitoring và backup.
    \item Bổ sung automated tests cho Phase 2 và unit tests cho service layer.
\end{itemize}

\section*{5. Kết luận chung}

Tổng kết lại, đồ án đã xây dựng được một hệ thống quản lý phòng mạch tư nhân có cấu trúc rõ ràng, có bằng chứng kỹ thuật, có quy trình phát triển theo phase và có khả năng mở rộng. Dù chưa đạt mức production-ready, sản phẩm phù hợp với mục tiêu học phần SE104 vì thể hiện đầy đủ các bước của quy trình kỹ nghệ phần mềm: đặc tả yêu cầu, thiết kế hệ thống, thiết kế phần mềm, hiện thực, kiểm thử, triển khai và đánh giá kết quả.
